\documentclass[11pt]{article}
\usepackage[margin=1in]{geometry}
\usepackage{amsmath,amssymb,amsthm}
\usepackage{hyperref}

\newtheorem{theorem}{Theorem}
\newtheorem{proposition}[theorem]{Proposition}

\title{A Finite Irregular $3$-Group Counterexample to Manna--Mandal Problem~7.10: Condition~(H1) Does Not Imply~(H2)}
\author{Salem Levinstein}
\date{November 26, 2026}

\begin{document}
\maketitle

\begin{abstract}
We give a negative answer to the second question of Problem~7.10 of Manna and Mandal (\texttt{arXiv:2611.27642v1}). For a nilpotent group $K$ of odd order and $G = K \times \mathbb{Z}_2$, Manna and Mandal established two necessary conditions (H1) and (H2) for the prime index elements graph $\Gamma(G)$ to be Hamiltonian, showed that (H1) implies (H2) whenever the Sylow subgroups of $K$ are abelian, and asked whether (H1) implies (H2) in general. We construct an irregular $3$-group $P$ of order $243$ such that $K = P \times \mathbb{Z}_5$ satisfies (H1) for every prime dividing $|K|$ ($N_3(P) = 4 \le 10$ and $N_5(\mathbb{Z}_5) = 1 \le 486$) while violating (H2) at the central subgroup $C = \langle u \rangle$ of order~$3$ ($c_P(C) = 30 > 19$). Directly, removing the $20$ vertices of $G = P \times \mathbb{Z}_{10}$ whose $3$-part generates~$C$ disconnects $\Gamma(G)$ into $31$ components, proving that $\Gamma(G)$ admits neither a Hamiltonian cycle nor a Hamiltonian path.
\end{abstract}

\section{Introduction}

Let $G$ be a finite group. The \emph{prime index elements graph} $\Gamma(G)$ has vertex set $G$, with distinct elements $x, y \in G$ adjacent if one of $\langle x \rangle, \langle y \rangle$ has prime index in the other~\cite{MM2026}. For a $p$-group $P$ and a cyclic subgroup $C \le P$, write $c_P(C)$ for the number of cyclic subgroups of $P$ containing $C$ as a maximal subgroup, and set $N_p(P) := c_P(1)$ (the number of subgroups of order~$p$).

Let $K$ be a nilpotent group of odd order and $G = K \times \mathbb{Z}_2$. For each prime $p \mid |K|$ with Sylow $p$-subgroup $P \le K$ and $m_p := |K|/|P|$, Manna and Mandal~\cite[Theorem~7.6]{MM2026} proved that Hamiltonicity of $\Gamma(G)$ requires both
\begin{align*}
  \text{(H1)} &\quad N_p(P) \le 2 m_p, \\
  \text{(H2)} &\quad c_P(C) \le 2\varphi(|C|)\, m_p - 1 \quad \text{for every cyclic subgroup } 1 \ne C \le P.
\end{align*}
They showed that (H1) implies (H2) when the Sylow subgroups of $K$ are abelian~\cite[Corollary~7.8(iii)]{MM2026}, and asked in~\cite[Problem~7.10]{MM2026}: \emph{``Does (H1) imply (H2), as it does for abelian Sylow subgroups?''} We show that the answer is \textbf{no}. (The first and third questions of Problem~7.10 remain open.)

\section{The Counterexample}

Let $A = \mathbb{Z}_9 \times \mathbb{Z}_9$ written additively, let $s : A \to A$ be the linear map $s(a, b) = (-b, a - b)$, and set $u = (-3, 3) = (6, 3) \in A$. Define $P = A \times \{0, 1, 2\}$ with multiplication
\[
  (v, i) \cdot (w, j) = \bigl(v + s^i(w) + \varepsilon_{i, j}\, u,\; (i + j) \bmod 3\bigr), \qquad
  \varepsilon_{i, j} := \begin{cases} 1 & \text{if } i + j \ge 3, \\ 0 & \text{otherwise.} \end{cases}
\]
Equivalently, $P = \langle a, b, x \mid a^9 = b^9 = [a, b] = 1,\; xax^{-1} = b,\; xbx^{-1} = a^{-1}b^{-1},\; x^3 = a^{-3}b^3 \rangle$.

\begin{theorem}\label{thm:main}
Let $K = P \times \mathbb{Z}_5$ and $G = K \times \mathbb{Z}_2 \cong P \times \mathbb{Z}_{10}$. Then:
\begin{enumerate}
  \item $P$ is a non-abelian $3$-group of order $243$, so $K$ is nilpotent of odd order $1215$.
  \item \textnormal{(H1)} holds for every prime dividing $|K|$: $N_3(P) = 4 \le 10 = 2 m_3$ and $N_5(\mathbb{Z}_5) = 1 \le 486 = 2 m_5$.
  \item \textnormal{(H2)} fails for $p = 3$ at $C = \langle (u, 0) \rangle$ of order~$3$: $c_P(C) = 30 > 19 = 2\varphi(3)m_3 - 1$.
  \item Deleting the $20$-vertex cut set $S = \{(y, t) \in P \times \mathbb{Z}_{10} : y \in \{(u, 0), (2u, 0)\}\}$ separates $\Gamma(G)$ into $31$ connected components; hence $\Gamma(G)$ has neither a Hamiltonian cycle nor a Hamiltonian path.
\end{enumerate}
\end{theorem}

\begin{proof}
\emph{(1) Group axioms.} On the standard basis of $A = \mathbb{Z}_9^2$, $s(1, 0) = (0, 1)$, $s^2(1, 0) = (-1, -1)$, $s^3(1, 0) = (1, 0)$ and similarly $s^3(0, 1) = (0, 1)$, so $s^3 = \mathrm{id}_A$ and $1 + s + s^2 = 0$. Also $s(6, 3) = (-3, 3) = (6, 3)$, so $s(u) = u$. For $(v, i), (w, j), (r, k) \in P$, the first coordinates of $((v, i)(w, j))(r, k)$ and $(v, i)((w, j)(r, k))$ are both
\[
  v + s^i(w) + s^{i + j}(r) + \lfloor (i + j + k) / 3 \rfloor\, u,
\]
and the second coordinates are both $(i + j + k) \bmod 3$, proving associativity. The element $(0, 0)$ is a two-sided identity, and $(v, i)$ has inverse $(-v, 0)$ for $i = 0$ and $(s^{3 - i}(-v - u), 3 - i)$ for $i \in \{1, 2\}$. Thus $|P| = 9^2 \cdot 3 = 243 = 3^5$.

\emph{(2)--(3) Orders, $N_3(P) = 4$, and $c_P(C) = 30$.} Because $1 + s + s^2 = 0$ and $s(u) = u$, every $v \in A$ satisfies
\[
  (v, 1)^3 = (v + s(v) + s^2(v) + u, 0) = (u, 0), \qquad (v, 2)^3 = (v + s^2(v) + s(v) + 2u, 0) = (2u, 0).
\]
Since $u = (6, 3)$ has order~$3$ in~$A$, all $162$ elements of $P \setminus (A \times \{0\})$ have order~$9$ and cube in $\{(u, 0), (2u, 0)\}$. Within $A \times \{0\} \cong \mathbb{Z}_9^2$, the elements of order~$3$ are the $8$ nonzero elements of $3A \times \{0\}$. Thus $P$ has exactly $8$ elements of order~$3$, yielding $N_3(P) = 8 / 2 = 4 \le 10 = 2m_3$. For $p = 5$, the Sylow $5$-subgroup $\mathbb{Z}_5$ has $m_5 = 243$ and $N_5(\mathbb{Z}_5) = 1 \le 486$.

A cyclic subgroup $D \le P$ contains $C = \langle (u, 0) \rangle$ as a maximal subgroup if and only if $D = \langle y \rangle$ with $o(y) = 9$ and $y^3 \in \{(u, 0), (2u, 0)\}$. Besides the $162$ elements with $i \in \{1, 2\}$, the elements $(v, 0) \in A \times \{0\}$ satisfying $3v \in \{u, 2u\}$ are the lifts of $(2, 1)$ and $(1, 2) \pmod 3$, giving $2 \cdot 3^2 = 18$ elements. Each cyclic subgroup of order~$9$ has $\varphi(9) = 6$ generators, so $c_P(C) = (162 + 18) / 6 = 30 > 19 = 2\varphi(3)m_3 - 1$.

\emph{(4) Direct non-Hamiltonicity.} Let $S = \{(y, t) \in P \times \mathbb{Z}_{10} : y \in \{(u, 0), (2u, 0)\}\}$, so $|S| = 20$. Define a $31$-valued map $\beta$ on $G \setminus S$ by $\beta(y, t) = \langle y \rangle$ if $y^3 \in \{(u, 0), (2u, 0)\}$ and $\beta(y, t) = \bot$ otherwise. Every one of the $30$ cyclic subgroups $D > C$ of order~$9$ is realized by $(y, 0)$ for a generator $y$ of $D$, and $\bot$ is realized by $(1_P, 0)$, so $|\beta(G \setminus S)| = 31$.

Suppose $g = (g_3, g')$ and $h = (h_3, h')$ in $G \setminus S$ are adjacent in $\Gamma(G)$; without loss of generality $\langle g \rangle < \langle h \rangle$ with prime index~$r$. Then $\langle g_3 \rangle \le \langle h_3 \rangle$ in $P$ with index dividing~$r$, hence equal to~$1$ or~$3$.
\begin{itemize}
  \item If $|\langle g_3 \rangle| = |\langle h_3 \rangle|$, then $\langle g_3 \rangle = \langle h_3 \rangle$, so $\beta(g) = \beta(h)$.
  \item If $|\langle h_3 \rangle| = 3|\langle g_3 \rangle|$, then $|\langle g_3 \rangle| \le 3$, so $g_3^3 = 1_P$ and $\beta(g) = \bot$. If $h_3^3 \in \{(u, 0), (2u, 0)\}$, then $|\langle h_3 \rangle| = 9$, forcing $|\langle g_3 \rangle| = 3$ and hence $\langle g_3 \rangle = C$, i.e., $g \in S$, a contradiction. Thus $\beta(h) = \bot = \beta(g)$.
\end{itemize}
Hence $\beta$ is constant on every edge of $\Gamma(G) - S$, so $\Gamma(G) - S$ has at least $31$ connected components (direct computation in \texttt{verify.py} confirms exactly~$31$). Deleting $20$ vertices from a Hamiltonian path leaves at most $21$ contiguous arcs, and deleting $20$ vertices from a Hamiltonian cycle leaves at most $20$. Since $31 > 21$, $\Gamma(G)$ has neither a Hamiltonian path nor a Hamiltonian cycle.
\end{proof}

\noindent\textbf{Remark.} By P.~Hall's theory of regular $p$-groups, any regular $p$-group satisfies the same $\Omega_1$-bound as abelian $p$-groups, so any counterexample $P$ must be irregular. Here $P$ is indeed irregular: $(u, 0)$ has $90$ cube roots while $|\Omega_1(P)| = 9$. Theorem~\ref{thm:main} is formally verified in Lean~4 (\texttt{MannaMandalPrimeIndexGraph.lean}) without \texttt{sorry} or non-standard axioms.

\begin{thebibliography}{9}
\bibitem{MM2026}
P.~Manna and S.~Mandal,
\emph{On the prime index elements graph of finite groups},
Preprint \texttt{arXiv:2611.27642v1} [math.GR], 20~November 2026.
\end{thebibliography}

\end{document}
